% GENERATED FILE. Edit canonical lessons, then run npm run generate:books.
% canonical-source-hash: fnv1a64:f7cd97f26857624b
% canonical-lessons: IT-C55-colazione, IT-C55-pera, IT-C55-riso, IT-C55-formaggio, IT-C55-uovo

\chapter{Breakfast, Pear, Rice, Cheese, Egg}
\label{ch:it-breakfast-pear-rice-cheese-egg}
\hlchaptermodality{55}

\begin{chapteropening}
\textbf{By the end of this chapter:} \emph{I can name breakfast, a pear, rice, cheese, an egg.}

Complete the last lesson of chapter 55: I can name breakfast, a pear, rice, cheese, an egg.
\end{chapteropening}

\section[la colazione]{la colazione — breakfast}

\label{lesson:IT-C55-colazione}



\begin{quote}
Before the new one: say the Italian for oil, then the Italian for lunch.
\end{quote}

\subsection*{You'll want to know: la colazione}
\textbf{la colazione} — \textquotedblleft{}breakfast\textquotedblright{}.

\begin{cousinweb}
\textbf{colazione} $\leftarrow$ Latin \emph{collatio}, \textquotedblleft{}a bringing together\textquotedblright{} — a learned borrowing. (EtymDB 2.1)
\end{cousinweb}

\begin{grammarlens}[title={What you've built}]
One more word for this chapter.
\end{grammarlens}

\subsection*{Your turn}
\begin{itemize}
\raggedright
  \item \emph{Say it:} \emph{la colazione}
  \item \emph{Say it:} \emph{la colazione}, once more
  \item \emph{Say it:} say \emph{la colazione}
  \item \emph{Recall:} say the Italian for oil, then the Italian for lunch, then say \emph{la colazione} again
\end{itemize}

\begin{tcolorbox}[breakable,colback=teal!4,colframe=teal!35!black,title={Before you move on}]
What does la colazione mean? (\textquotedblleft{}Breakfast\textquotedblright{}.) Say it once more.
\end{tcolorbox}
\section[la pera]{la pera — a pear}

\label{lesson:IT-C55-pera}



\begin{quote}
Before the new one: say the Italian for lunch, then the Italian for breakfast.
\end{quote}

\subsection*{You'll want to know: la pera}
\textbf{la pera} — \textquotedblleft{}a pear\textquotedblright{}.

\begin{cousinweb}
\textbf{pera} $\leftarrow$ Vulgar Latin \emph{pira}, from Latin \emph{pirum}, \textquotedblleft{}pear\textquotedblright{} — inherited. (EtymDB 2.1)
\end{cousinweb}

\begin{grammarlens}[title={What you've built}]
One more word for this chapter.
\end{grammarlens}

\subsection*{Your turn}
\begin{itemize}
\raggedright
  \item \emph{Say it:} \emph{la pera}
  \item \emph{Say it:} \emph{la pera}, once more
  \item \emph{Say it:} say \emph{la pera}
  \item \emph{Recall:} say the Italian for lunch, then the Italian for breakfast, then say \emph{la pera} again
\end{itemize}

\begin{tcolorbox}[breakable,colback=teal!4,colframe=teal!35!black,title={Before you move on}]
What does la pera mean? (\textquotedblleft{}A pear\textquotedblright{}.) Say it once more.
\end{tcolorbox}
\section[il riso]{il riso — rice}

\label{lesson:IT-C55-riso}



\begin{quote}
Before the new one: say the Italian for breakfast, then the Italian for a pear.
\end{quote}

\subsection*{You'll want to know: il riso}
\textbf{il riso} — \textquotedblleft{}rice\textquotedblright{}.

\begin{cousinweb}
\textbf{riso} $\leftarrow$ Latin \emph{oryza}, \textquotedblleft{}rice\textquotedblright{} — inherited; Latin took it from Greek \emph{oryza}. (EtymDB 2.1)
\end{cousinweb}

\begin{grammarlens}[title={What you've built}]
One more word for this chapter.
\end{grammarlens}

\subsection*{Your turn}
\begin{itemize}
\raggedright
  \item \emph{Say it:} \emph{il riso}
  \item \emph{Say it:} \emph{il riso}, once more
  \item \emph{Say it:} say \emph{il riso}
  \item \emph{Recall:} say the Italian for breakfast, then the Italian for a pear, then say \emph{il riso} again
\end{itemize}

\begin{tcolorbox}[breakable,colback=teal!4,colframe=teal!35!black,title={Before you move on}]
What does il riso mean? (\textquotedblleft{}Rice\textquotedblright{}.) Say it once more.
\end{tcolorbox}
\section[il formaggio]{il formaggio — cheese}

\label{lesson:IT-C55-formaggio}



\begin{quote}
Before the new one: say the Italian for a pear, then the Italian for rice.
\end{quote}

\subsection*{You'll want to know: il formaggio}
\textbf{il formaggio} — \textquotedblleft{}cheese\textquotedblright{}.

\begin{cousinweb}
\textbf{formaggio} $\leftarrow$ Old French \emph{formage} — borrowed; behind it is Latin \emph{formaticum}, cheese shaped in a mould. (EtymDB 2.1)
\end{cousinweb}

\begin{grammarlens}[title={What you've built}]
One more word for this chapter.
\end{grammarlens}

\subsection*{Your turn}
\begin{itemize}
\raggedright
  \item \emph{Say it:} \emph{il formaggio}
  \item \emph{Say it:} \emph{il formaggio}, once more
  \item \emph{Say it:} say \emph{il formaggio}
  \item \emph{Recall:} say the Italian for a pear, then the Italian for rice, then say \emph{il formaggio} again
\end{itemize}

\begin{tcolorbox}[breakable,colback=teal!4,colframe=teal!35!black,title={Before you move on}]
What does il formaggio mean? (\textquotedblleft{}Cheese\textquotedblright{}.) Say it once more.
\end{tcolorbox}
\section[l'uovo]{l'uovo — an egg}

\label{lesson:IT-C55-uovo}



\begin{quote}
Before the new one: say the Italian for rice, then the Italian for cheese.
\end{quote}

\subsection*{You'll want to know: l'uovo}
\textbf{l'uovo} — \textquotedblleft{}an egg\textquotedblright{}.

\begin{cousinweb}
\textbf{uovo} $\leftarrow$ Latin \emph{ovum}, \textquotedblleft{}egg\textquotedblright{} — inherited. (EtymDB 2.1)
\end{cousinweb}

\begin{grammarlens}[title={What you've built}]
One more word for this chapter.
\end{grammarlens}

\subsection*{Your turn}
\begin{itemize}
\raggedright
  \item \emph{Say it:} \emph{l'uovo}
  \item \emph{Say it:} \emph{l'uovo}, once more
  \item \emph{Say it:} say \emph{l'uovo}
  \item \emph{Recall:} say the Italian for rice, then the Italian for cheese, then say \emph{l'uovo} again
\end{itemize}

\begin{tcolorbox}[breakable,colback=teal!4,colframe=teal!35!black,title={Before you move on}]
What does l'uovo mean? (\textquotedblleft{}An egg\textquotedblright{}.) Say it once more.
\end{tcolorbox}
